% Vietnamese translation for OI Public Library
% Source-PDF: IOI中国国家候选队论文/2002/俞玮.pdf
\documentclass[11pt]{article}
\usepackage{oipl-vn}
\usepackage{tikz}
\usetikzlibrary{arrows.meta,patterns}

\title{Trò chơi Ulam và mã hóa}
\author{Yu Wei}
\date{}

\definecolor{markgreen}{RGB}{143,207,0}
\definecolor{lightblue}{RGB}{198,230,232}
\definecolor{softgreen}{RGB}{178,205,91}
\definecolor{badred}{RGB}{230,35,30}

\newcommand{\yes}{\mathrm{Y}}
\newcommand{\no}{\mathrm{N}}
\newcommand{\xor}{\oplus}

\begin{document}
\maketitle

\origtitle{Ulam's game and coding}
\sourcepath{IOI China candidate team papers / 2002 / Yu Wei.pdf}

\section{Trò chơi Ulam}

Trong trò chơi Ulam có \(n\) số. Trong số đó có đúng một số khác với các số
còn lại; đó là số mà người chơi đã chọn. Ta được đặt các câu hỏi dạng
``số ấy có nằm trong tập \(X\) hay không?'' và mỗi câu hỏi chỉ nhận được câu
trả lời ``có'' hoặc ``không''. Khó khăn nằm ở chỗ trong toàn bộ quá trình có
thể xuất hiện một câu trả lời sai, nhưng số lần sai không vượt quá một.

Mục tiêu là tìm ra số đã được chọn bằng càng ít câu hỏi càng tốt.

\section{Bài toán gốc không có sai sót}

Trước hết xét phiên bản đơn giản: vẫn có \(n\) số và một số được chọn, cách
hỏi vẫn là hỏi số được chọn có thuộc một tập \(X\) hay không, nhưng mọi câu
trả lời đều đúng.

Lời giải quen thuộc là \term{binary search}. Mỗi lần ta chia tập các số còn
có thể là đáp án thành hai phần có kích thước gần bằng nhau nhất, lấy một
phần làm tập \(X\), rồi dựa vào câu trả lời để loại bỏ phần không thể chứa
đáp án. Sau khoảng \(\lg n\) câu hỏi chỉ còn lại một số, và số đó là đáp án.

Nhìn theo một cách khác, binary search chỉ là một phương pháp sàng lọc hơi
tham lam: ở mỗi bước nó cố gắng loại bỏ nhiều số nhất trong trường hợp xấu
nhất.

\begin{center}
\begin{tikzpicture}[x=0.75cm,y=0.42cm]
  \tikzset{
    cell/.style={draw=black!35,minimum width=0.75cm,minimum height=0.42cm,inner sep=0pt},
    arr/.style={-{Latex[length=3mm]},line width=0.6pt,draw=black!65}
  }
  \foreach \stage/\x/\top/\bottom in {0/0/8/0,1/3.2/4/0,2/6.4/2/0,3/9.6/1/0} {
    \foreach \r in {1,...,8} {
      \pgfmathtruncatemacro{\keep}{\r>\bottom && \r<=\top}
      \ifnum\keep=1
        \node[cell,fill=markgreen] at (\x,\r) {};
      \else
        \node[cell,fill=lightblue] at (\x,\r) {};
      \fi
    }
    \node[below=0.25cm] at (\x,0.5) {bước \(\stage\)};
  }
  \draw[arr] (0.6,4.5) -- (2.6,4.5);
  \draw[arr] (3.8,4.5) -- (5.8,4.5);
  \draw[arr] (7.0,4.5) -- (9.0,4.5);
\end{tikzpicture}
\end{center}

Sơ đồ trên chỉ biểu diễn ý tưởng sàng lọc: phần xanh là tập còn có thể chứa
đáp án, và sau mỗi câu hỏi ta giữ lại phần phù hợp với câu trả lời.

\section{Binary search khi có một câu trả lời sai}

Khi có thể có một câu trả lời sai, việc chia đôi theo số lượng không còn đủ.
Nguyên tắc đúng hơn là chia đều \emph{khả năng} chứa số được chọn: hai phần
sau khi hỏi phải có xác suất chứa đáp án như nhau, xét theo mọi kịch bản sai
sót còn hợp lệ.

Một cách hiện thực là phân loại các số theo số lần câu trả lời trước đó phải
sai thì số ấy mới còn có thể là đáp án. Với giả thiết nhiều nhất một lỗi, ta
có hai lớp:
\begin{itemize}
  \item lớp \(0\): các số còn phù hợp nếu chưa có câu trả lời sai nào;
  \item lớp \(1\): các số chỉ còn phù hợp nếu đã có đúng một câu trả lời sai.
\end{itemize}
Trong cùng một lớp, các số có cùng mức khả dĩ. Vì vậy ở mỗi câu hỏi ta cố
gắng chia đều từng lớp, lấy một phần của mỗi lớp đưa vào tập \(X\).

\subsection{Một thao tác mẫu}

Hình sau mô phỏng bốn trạng thái liên tiếp. Mỗi khối gồm hai cột: cột trái là
lớp \(0\), cột phải là lớp \(1\). Các ô xanh là những số còn hợp lệ trong lớp
tương ứng; phần gạch chéo là phần bị xác định là không chứa số được chọn sau
câu trả lời hiện tại và sẽ chuyển sang lớp lỗi cao hơn hoặc bị loại.

\begin{center}
\begin{tikzpicture}[x=0.52cm,y=0.44cm]
  \tikzset{
    box/.style={draw=black,minimum width=0.52cm,minimum height=0.44cm,inner sep=0pt},
    arrow/.style={-{Latex[length=3mm]},line width=0.6pt}
  }
  \newcommand{\stage}[4]{
    \begin{scope}[shift={(#1,0)}]
      \foreach \r in {1,...,8} {
        \node[box] at (0,\r) {};
        \node[box] at (1,\r) {};
      }
      #2
      #3
      \node[below] at (0.5,0.35) {#4};
    \end{scope}
  }
  \stage{0}{
    \foreach \r in {1,...,8} \node[box,fill=markgreen] at (0,\r) {};
  }{
    \foreach \r in {1,...,4} \node[box,pattern=north east lines,pattern color=markgreen!80!black] at (0,\r) {};
  }{trước hỏi}
  \stage{4.2}{
    \foreach \r in {1,2,5,6} \node[box,fill=markgreen] at (0,\r) {};
    \foreach \r in {3,4} \node[box,fill=markgreen] at (1,\r) {};
  }{
    \foreach \r in {5,6} \node[box,pattern=north east lines,pattern color=markgreen!80!black] at (0,\r) {};
  }{sau hỏi 1}
  \stage{8.4}{
    \foreach \r in {2} \node[box,fill=markgreen] at (0,\r) {};
    \foreach \r in {3,4,7} \node[box,fill=markgreen] at (1,\r) {};
  }{
    \foreach \r in {5,6,8} \node[box,pattern=north east lines,pattern color=markgreen!80!black] at (1,\r) {};
  }{sau hỏi 2}
  \stage{12.6}{
    \foreach \r in {2} \node[box,fill=markgreen] at (0,\r) {};
    \foreach \r in {3,4,7,8} \node[box,fill=markgreen] at (1,\r) {};
  }{
  }{sau hỏi 3}
  \draw[arrow] (1.7,4.5) -- (3.3,4.5);
  \draw[arrow] (5.9,4.5) -- (7.5,4.5);
  \draw[arrow] (10.1,4.5) -- (11.7,4.5);
\end{tikzpicture}
\end{center}

Điểm quan trọng của hình không phải là vị trí cụ thể của từng ô, mà là cách
mỗi câu hỏi chia đều từng lớp lỗi rồi cập nhật lớp của các số theo câu trả lời.

\subsection{Ước lượng số câu hỏi}

Giả sử dùng cách sàng lọc trên. Sau \(k\) câu hỏi, số phần tử còn trong lớp
\(0\) là
\[
  \frac{n}{2^k}.
\]
Gọi \(f(k)\) là số phần tử còn trong lớp \(1\). Mỗi bước giữ lại một nửa lớp
\(1\) cũ, đồng thời một nửa phần không khớp của lớp \(0\) chuyển sang lớp
\(1\). Do đó
\[
  f(k)=\frac{f(k-1)}{2}+\frac{n}{2^k}.
\]
Từ truy hồi này suy ra
\[
  f(k)=\frac{kn}{2^k}.
\]

Khi tổng số phần tử ở hai lớp không vượt quá \(1\), đáp án đã được xác định:
\[
  f(k)+\frac{n}{2^k}\le 1.
\]
Thay công thức của \(f(k)\), ta được điều kiện
\[
  2^k \ge kn+n.
\]
Vì vậy số câu hỏi của phương pháp sàng lọc là nghiệm nguyên nhỏ nhất \(k\)
thỏa bất đẳng thức trên.

\section{Phương pháp mới: mã hóa}

Phương pháp sàng lọc có ba điểm bất tiện. Nó phải lưu trạng thái trung gian,
phụ thuộc trực tiếp vào kết quả của bước trước, và cần xử lý cẩn thận các
trường hợp biên.

Ý tưởng mã hóa tránh các bất tiện ấy bằng cách gán trước cho mỗi số một mã
nhị phân cố định. Câu hỏi thứ \(i\) sẽ hỏi tập các số có bit thứ \(i\) của mã
bằng \(1\). Câu trả lời được đổi thành mã kết quả theo quy ước
\[
  \yes \leftrightarrow 1,\qquad \no \leftrightarrow 0.
\]
Nếu không có lỗi, đáp án chính là số có mã trùng với mã kết quả.

\subsection{Ví dụ mã hóa}

Với \(8\) số, dùng ba bit thông tin là đủ. Gán mã theo bảng sau:
\[
\begin{array}{c|ccc}
\text{số} & X_1 & X_2 & X_3\\ \hline
1 & 0 & 0 & 0\\
2 & 0 & 0 & 1\\
3 & 0 & 1 & 0\\
4 & 0 & 1 & 1\\
5 & 1 & 0 & 0\\
6 & 1 & 0 & 1\\
7 & 1 & 1 & 0\\
8 & 1 & 1 & 1
\end{array}
\]

Nếu ba câu trả lời là \(\yes,\yes,\no\), mã kết quả là \(110\). Khi đó hàng
tương ứng với số \(7\) được đánh dấu ở cả ba cột phù hợp, nên số \(7\) là đáp
án.

\begin{center}
\begin{tikzpicture}[x=0.55cm,y=0.45cm]
  \tikzset{box/.style={draw=black,minimum width=0.55cm,minimum height=0.45cm,inner sep=0pt}}
  \foreach \c/\lab in {1/\(X_1\),2/\(X_2\),3/\(X_3\)} {
    \node at (\c,9) {\lab};
  }
  \foreach \r in {1,...,8} {
    \node[left] at (0.55,9-\r) {\r};
    \foreach \c in {1,2,3} {
      \node[box] at (\c,9-\r) {};
    }
  }
  \foreach \r/\c in {5/1,6/1,7/1,8/1,3/2,4/2,7/2,8/2,2/3,4/3,6/3,8/3} {
    \node[box,fill=markgreen] at (\c,9-\r) {};
  }
  \draw[badred,line width=1pt] (0.7,2) rectangle (3.3,2.45);
  \node[right,badred] at (3.65,2.22) {\(110\)};
  \node[below] at (2,0.15) {hàng \(7\) trùng mã kết quả};
\end{tikzpicture}
\end{center}

\section{Lỗi và mã sửa lỗi}

Khi trong các câu trả lời có một lỗi, mã kết quả nhận được cũng sai đúng một
bit. Vì thế các khả năng trả lời khác nhau có thể được xem như các mã lỗi
khác nhau. Ta cần thiết kế \term{mã sửa lỗi} sao cho từ mã kết quả bị sai
một bit vẫn suy ra được mã đúng.

Một cách tự nhiên là dùng mã kiểm tra chẵn lẻ. Các bit kiểm tra được định
nghĩa là tổng nhị phân, tức phép xor, của một số bit thông tin. Khi một bit
bị sai, mọi đẳng thức kiểm tra liên quan đến bit ấy sẽ không còn đúng.

Giả sử tạm thời các bit không phải bit kiểm tra đều đúng. Ta xét những bit
kiểm tra bị ``sai'', tức những phương trình kiểm tra không thỏa:
\begin{itemize}
  \item nếu không có phương trình nào sai, thì không có lỗi;
  \item nếu chỉ có một phương trình sai, lỗi nằm ở chính bit kiểm tra ấy;
  \item nếu có nhiều phương trình sai, lỗi nằm ở bit thông tin xuất hiện
  đúng trong các phương trình kiểm tra ấy.
\end{itemize}

\subsection{Cấu tạo bit kiểm tra}

Giả sử có \(m=4\) bit thông tin \(X_1,X_2,X_3,X_4\). Ta cần thêm \(r\) bit
kiểm tra sao cho mỗi vị trí bit, cùng với trường hợp không có lỗi, có một
dấu hiệu phân biệt. Vì vậy phải có
\[
  2^r \ge m+r+1.
\]
Với \(m=4\), giá trị nhỏ nhất là \(r=3\). Đặt ba bit kiểm tra là
\(X_5,X_6,X_7\) và chọn các phương trình
\[
\begin{aligned}
  X_5 &= X_2 \xor X_3 \xor X_4,\\
  X_6 &= X_1 \xor X_3 \xor X_4,\\
  X_7 &= X_1 \xor X_2 \xor X_4.
\end{aligned}
\]

Mỗi bit thông tin xuất hiện trong một tập phương trình khác nhau:
\[
\begin{array}{c|ccc}
 & X_5 & X_6 & X_7\\ \hline
X_1 & 0 & 1 & 1\\
X_2 & 1 & 0 & 1\\
X_3 & 1 & 1 & 0\\
X_4 & 1 & 1 & 1
\end{array}
\]
Các cột riêng \(100,010,001\) lần lượt nhận diện lỗi ở \(X_5,X_6,X_7\), còn
\(000\) biểu thị không có lỗi. Như vậy mọi lỗi một bit đều có một dấu hiệu
riêng.

\subsection{Đối ứng với câu hỏi}

Phép xor trên mã tương đương với phép xor trên chính các câu hỏi. Chẳng hạn
từ
\[
  X_5 = X_2 \xor X_3 \xor X_4
\]
ta hiểu rằng câu hỏi thứ \(5\) hỏi tập đối xứng sai khác của ba tập trong
các câu hỏi \(2,3,4\). Một số thuộc \(X_5\) khi và chỉ khi nó thuộc lẻ số
trong ba tập \(X_2,X_3,X_4\).

Vì vậy việc thêm bit kiểm tra vào mã không chỉ là thao tác trên ký hiệu nhị
phân; nó còn cho ta cách tạo thêm các câu hỏi kiểm tra trong trò chơi Ulam.

\section{Số câu hỏi}

Kết quả đếm câu hỏi của phương pháp sàng lọc đã thu được ở trên: số câu hỏi
là nghiệm nguyên nhỏ nhất \(k\) thỏa
\[
  2^k \ge kn+n.
\]

Với phương pháp mã hóa, đặt \(m=\lg n\) là số bit thông tin cần có khi không
có lỗi, và đặt \(t\) là tổng số câu hỏi sau khi thêm bit kiểm tra. Ngoài
\(m\) bit thông tin, mỗi khả năng lỗi phải ứng với một mã kiểm tra riêng.
Do đó cần
\[
  2^{t-m}\ge t+1.
\]
Mặt khác, tất cả khả năng gồm giá trị thật và vị trí lỗi cũng phải có mã
khác nhau, nên điều kiện tương đương là
\[
  2^t \ge 2^m(t+1).
\]
Kết quả là \(t\) là nghiệm nguyên nhỏ nhất thỏa
\[
  2^{t-m}\ge t+1.
\]

Hai công thức trùng nhau khi \(m=\lg n\):
\[
  2^t\ge tn+n
  \quad\Longleftrightarrow\quad
  2^t\ge t2^m+2^m
  \quad\Longleftrightarrow\quad
  2^{t-m}\ge t+1.
\]
Nói cách khác, theo số câu hỏi trong trường hợp một lỗi, cách sàng lọc và
cách mã hóa đạt cùng ngưỡng.

\section{Mở rộng}

\subsection{Nhiều lỗi hơn}

Khi số câu trả lời sai có thể lớn hơn một, việc mở rộng trực tiếp phương pháp
mã hóa trở nên khó hơn nhiều, vì mã phải sửa được nhiều lỗi. Với hướng sàng
lọc, ý tưởng lại khá trực tiếp: nếu có nhiều nhất \(e\) lỗi, ta chia các số
thành \(e+1\) lớp, trong đó lớp \(i\) gồm các số còn hợp lệ nếu đã có đúng
\(i\) câu trả lời sai. Mỗi lần hỏi, ta cố gắng chia đều từng lớp.

\begin{center}
\begin{tikzpicture}[x=0.5cm,y=0.42cm]
  \tikzset{box/.style={draw=black!65,minimum width=0.5cm,minimum height=0.42cm,inner sep=0pt}}
  \foreach \c in {0,...,3} {
    \node[above] at (\c,8.8) {lớp \(\c\)};
    \foreach \r in {1,...,8} \node[box] at (\c,9-\r) {};
  }
  \foreach \r in {1,2,3,4} \node[box,fill=markgreen] at (0,9-\r) {};
  \foreach \r in {5,6} \node[box,fill=markgreen!70] at (1,9-\r) {};
  \foreach \r in {7} \node[box,fill=softgreen] at (2,9-\r) {};
  \foreach \r in {8} \node[box,fill=softgreen!50] at (3,9-\r) {};
  \node[below] at (1.5,0.25) {mỗi câu hỏi chia từng lớp lỗi};
\end{tikzpicture}
\end{center}

\subsection{Nhiều loại câu trả lời và bài toán cân bi}

Nếu mỗi câu hỏi có nhiều hơn hai loại trả lời, phương pháp sàng lọc phải lưu
nhiều nhóm trạng thái hơn. Còn với phương pháp mã hóa, ta thay mã nhị phân
bằng mã trong cơ số lớn hơn và thiết kế mã sửa lỗi tương ứng.

Bài toán cân bi là ví dụ điển hình. Mỗi lần cân có ba kết quả: đĩa trái
nặng hơn, cân bằng, hoặc đĩa phải nặng hơn. Do đó biểu diễn tự nhiên là mã
tam phân. Ngoài giá trị của viên bi đặc biệt, kiểu sai lệch cũng cần được mã
hóa, rồi ta thiết kế mã sửa lỗi tam phân để nhận diện lỗi.

\section{Kết luận}

Trò chơi Ulam cho thấy cùng một bài toán hỏi đáp có thể được nhìn theo hai
cách bổ sung nhau. Cách sàng lọc quản lý trực tiếp tập các ứng viên và các
lớp lỗi; nó mở rộng tự nhiên sang nhiều lỗi bằng cách thêm lớp. Cách mã hóa
gán trước biểu diễn cho từng số; trong trường hợp một lỗi, mã kiểm tra chẵn
lẻ cho phép xác định vị trí lỗi và đạt cùng số câu hỏi tối ưu như phân tích
sàng lọc.

Các hướng mở rộng chính là sàng lọc khi có nhiều lỗi, lựa chọn biểu diễn mã,
thiết kế mã sửa lỗi, và các hệ mã nhiều cơ số cho bài toán có nhiều loại câu
trả lời.

\end{document}
